\section{Exploratory Data Analysis}

The exploratory analysis aims to identify which signals are reliable, redundant, or structurally informative for the classification task, with the goal of motivating preprocessing and modeling choices.

\paragraph{Temporal Metadata and Redundancy}
The \textit{timestamp} attribute contains missing values (approximately \textbf{[TIMESTAMP\_MISSING\_PERCENT]\%}) and exhibits a coarse temporal granularity, with articles concentrated in a limited number of publication years. Finer resolutions (month or day) do not show deterministic relationships with the target labels.

An auxiliary experiment demonstrates that publication year can be predicted from the raw article text---including HTML structure---with accuracy well above chance (\textbf{[YEAR\_PRED\_ACCURACY]}), indicating that temporal information is implicitly encoded in linguistic and editorial patterns. As a result, the timestamp is considered a redundant signal and is excluded as an explicit feature.

\textbf{[FIGURE 1 HERE: Year distribution and year prediction from text]}

\paragraph{Non-informative Metadata}
The \textit{page\_rank} feature displays an approximately uniform distribution and does not exhibit a meaningful correlation with the target labels. Empirical checks confirm that its inclusion does not improve classification performance, and the feature is therefore discarded.

\paragraph{Structural Ambiguity}
The dataset contains duplicated or republished articles with identical content assigned to different labels (\textbf{[DUPLICATE\_RATE]\%}), introducing an intrinsic ambiguity that defines an upper bound on achievable performance. This motivates the use of macro-averaged F1 score as the primary evaluation metric.

\paragraph{Editorial Bias and Source Informativeness}
A strong dependency is observed between article source and target labels, reflecting editorial specialization across publishers. Encoding the source via one-hot representations yields a strong linear baseline and motivates a two-stage strategy that separates near-deterministic cases from semantically ambiguous ones.

\textbf{[FIGURE 2 HERE: Source $\times$ Label heatmap]}

\paragraph{Implications}
Overall, the dataset combines noisy HTML-heavy text, incomplete metadata, and strong editorial regularities. These properties motivate a progressive modeling approach that exploits structural signals early and reserves more expressive models for ambiguous samples.
